%! AP Statistics — Unit 4, Lesson 4.10 — Lesson Plan
\documentclass[10pt]{article}
\usepackage{apstats-article}
\usepackage{apstats-boxes}
\usepackage{graphicx}
\graphicspath{{images/}}

\newcommand{\UnitNumberName}{Unit 4: Probability, Random Variables, and Probability Distributions \quad}
\newcommand{\LessonNumberName}{Lesson 4.10: Introduction to the Binomial Distribution}

\begin{document}
\begin{center}
    {\Large\bfseries \CourseName: \SchoolYear \\
    {\small\bfseries \UnitNumberName \LessonNumberName}}
\end{center}

\begin{tcolorbox}[colback=sky, colframe=navy, boxrule=0.9pt, arc=2mm,
    left=2mm, right=2mm, top=1.5mm, bottom=1.5mm]
    \textbf{Primary Objective:} Students will determine whether a random variable follows
    a binomial distribution by checking the BINS conditions, then compute a single
    binomial probability $P(X = k)$ using the formula and calculator (Big Idea: UNC-3).
\end{tcolorbox}

\renewcommand{\arraystretch}{1.6}
\begin{skillbox}[Priority Ideas \& Skills]{goldbox}
\begin{minipage}[t]{0.48\linewidth}
    \textbf{AP Skill 3.A --- Determining Distributions}
    \begin{itemize}
        \item Check the four BINS conditions in context: Binary, Independent, fixed
              $N$ trials, Same probability $p$ (UNC-3.A).
        \item If all four conditions are met, state $X \sim B(n, p)$.
        \item Set up and compute a single probability $P(X = k)$ using the binomial
              formula and \texttt{binompdf(n, p, k)}.
    \end{itemize}
\end{minipage}
\hfill
\begin{minipage}[t]{0.48\linewidth}
    \textbf{Key Understandings}
    \begin{itemize}
        \item Each BINS condition must be explicitly verified from the problem context.
        \item The $10\%$ condition allows independence to be assumed when sampling
              without replacement from a large population.
        \item $P(X = k) = \binom{n}{k} p^k (1-p)^{n-k}$; the combination
              $\binom{n}{k}$ counts the number of ways to arrange $k$ successes
              among $n$ trials.
        \item ``Binomial'' refers to the two outcomes: success and failure.
    \end{itemize}
\end{minipage}
\end{skillbox}

\begin{skillbox}[{Vocabulary, Concepts, \& Theorems}]{greenbox}
\begin{minipage}[t]{0.48\linewidth}
    \begin{itemize}
        \item \textbf{Binomial setting} --- four conditions (BINS) are satisfied
        \item \textbf{Binomial random variable} --- counts successes in $n$ trials
        \item \textbf{BINS conditions} --- Binary, Independent, fixed $N$, Same $p$
        \item \textbf{$10\%$ condition} --- $n \le 0.10 \cdot N$ allows approximate
              independence
    \end{itemize}
\end{minipage}
\hfill
\begin{minipage}[t]{0.48\linewidth}
    \begin{itemize}
        \item \textbf{Binomial coefficient} $\binom{n}{k} = \dfrac{n!}{k!(n-k)!}$
        \item \textbf{Binomial probability formula} ---
              $P(X=k) = \binom{n}{k} p^k (1-p)^{n-k}$
        \item \textbf{binompdf(n, p, k)} --- calculator command for $P(X = k)$
        \item \textbf{Success / Failure} --- the two outcomes in a binary trial
    \end{itemize}
\end{minipage}
\end{skillbox}

\begin{skillbox}[Activate Prior Knowledge \& Spiral Review]{sky}
\begin{minipage}[t]{0.48\linewidth}
    \textbf{Prior Knowledge Check (3--4 min)}
    \begin{itemize}
        \item Quick review: What is a discrete random variable? (Lesson 4.5)
        \item Connect: We described probability distributions with tables; now we
              will use a formula to generate the whole distribution.
        \item Prompt: \textit{``A basketball player makes 70\% of free throws.
              If she shoots 5 in a row, list two things you'd need to assume to
              model this situation simply.''}
    \end{itemize}
\end{minipage}
\hfill
\begin{minipage}[t]{0.48\linewidth}
    \textbf{Connections Forward}
    \begin{itemize}
        \item Binomial mean and SD formulas $\to$ Lesson 4.11
        \item Cumulative binomial probabilities $\to$ Lesson 4.11
        \item Geometric distribution (trials until first success) $\to$ Lesson 4.12
        \item Normal approximation to binomial $\to$ Unit 5
    \end{itemize}
\end{minipage}
\end{skillbox}

\begin{skillbox}[Hook]{sky}
    \textbf{Hook (4--5 min)}\\
    Display: \textit{``A multiple-choice quiz has 10 questions, each with 4 options.
    A student guesses randomly on every question. What is the probability of getting
    exactly 6 correct?''} \\[4pt]
    Let students estimate. Then reveal that this probability is about 1.6\% ---
    surprisingly low given that the expected number correct is 2.5 out of 10.
    Ask: \textit{``What information did you need to set up this problem?''} Draw out:
    the number of trials, probability of success on each trial, and what counts as
    a ``success.'' Transition: today we formalize this into the binomial distribution.
\end{skillbox}

\begin{skillbox}[Lesson]{sky}
\begin{minipage}[t]{0.48\linewidth}
    \textbf{Part 1 --- BINS Conditions (8--10 min)}
    \begin{enumerate}
        \item Introduce the acronym BINS. For each letter, give the
              precise condition and a student-friendly checkpoint.
        \item Work a ``yes'' example and a ``no'' example side by side.
        \item Key point: sampling without replacement violates independence ---
              but the $10\%$ condition allows us to proceed when $n \le 0.10N$.
        \item Have students check BINS for the quiz context.
    \end{enumerate}
\end{minipage}
\hfill
\begin{minipage}[t]{0.48\linewidth}
    \textbf{Part 2 --- Binomial Formula (8--10 min)}
    \begin{enumerate}
        \item Write $P(X = k) = \binom{n}{k} p^k (1-p)^{n-k}$.
              Unpack each factor: the combination counts arrangements; $p^k$
              is the probability of $k$ successes in a row; $(1-p)^{n-k}$
              is the probability of $n-k$ failures.
        \item Model: solve the quiz problem step by step ($n=10$, $p=0.25$, $k=6$).
        \item Connect to calculator: \texttt{binompdf(10, 0.25, 6)}.
        \item Emphasize: on the AP exam, show the formula setup even if you
              use the calculator.
    \end{enumerate}
\end{minipage}
\end{skillbox}

\begin{skillbox}[Active Monitoring]{redbox}
\begin{minipage}[t]{0.48\linewidth}
    \textbf{Circulate \& Check For\ldots}
    \begin{itemize}
        \item Students who skip the BINS check and jump to the formula ---
              require written justification of each condition.
        \item Errors in evaluating $\binom{n}{k}$ (common: dividing $n!$ by $k!$
              without also dividing by $(n-k)!$).
        \item Confusing $p$ with $1-p$ (e.g., computing $P(\text{fail})^k$ instead
              of $P(\text{success})^k$).
    \end{itemize}
\end{minipage}
\hfill
\begin{minipage}[t]{0.48\linewidth}
    \textbf{Cold-Call Prompts}
    \begin{itemize}
        \item ``Walk me through why the I condition is satisfied here.''
        \item ``What does the $\binom{n}{k}$ term count in this problem?''
        \item ``What would change in your formula if I asked for exactly 3
              successes instead of exactly 6?''
    \end{itemize}
\end{minipage}
\end{skillbox}

\begin{skillbox}[Group Work \& Differentiation]{redbox}
\begin{minipage}[t]{0.48\linewidth}
    \textbf{Three-Tier Activity (10--12 min)}
    \begin{itemize}
        \item \textbf{Tier R}: Given two scenarios, students identify which is/is not
              a binomial setting and state which condition fails.
        \item \textbf{Tier A}: Two scenarios that satisfy BINS; students compute
              $P(X = k)$ using the formula and verify with \texttt{binompdf}.
        \item \textbf{Tier E}: A context where the $10\%$ condition is the only way
              independence can be justified; students argue in writing.
    \end{itemize}
\end{minipage}
\hfill
\begin{minipage}[t]{0.48\linewidth}
    \textbf{Differentiation}
    \begin{itemize}
        \item \textit{Support}: Provide a BINS checklist card with a sentence starter
              for each condition (``The outcomes are binary because\ldots'').
        \item \textit{Extension}: Challenge students to build the full probability
              distribution table for $n = 5$, $p = 0.3$ and verify that all
              probabilities sum to 1.
        \item \textit{ELL}: Post the BINS acronym with a visual for each letter.
    \end{itemize}
\end{minipage}
\end{skillbox}

\begin{skillbox}[Individual Work \& Assessment]{redbox}
\begin{minipage}[t]{0.48\linewidth}
    \textbf{Exit Ticket (5 min)}\\
    Scenario: \textit{``A factory produces light bulbs that are defective 8\% of the
    time. An inspector randomly selects 15 bulbs.''} Students (1) verify BINS, (2)
    define $X$, state its distribution, and (3) compute $P(X = 2)$.
\end{minipage}
\hfill
\begin{minipage}[t]{0.48\linewidth}
    \textbf{AP-Style Check}\\
    \textit{``Which scenario does NOT satisfy all conditions for a binomial
    distribution?''}
    \begin{enumerate}[label=(\Alph*)]
        \item Flipping a fair coin 20 times; $X =$ number of heads.
        \item Drawing 5 cards from a deck without replacement; $X =$ number of aces.
        \item Rolling a die 12 times; $X =$ number of 3s.
        \item Asking 50 randomly selected voters if they approve of a law; sample
              is less than 10\% of registered voters.
    \end{enumerate}
    Answer: (B) --- sampling without replacement and $n \ge 10\%$ of the deck.
\end{minipage}
\end{skillbox}

\begin{skillbox}[Reinforcement \& Extension]{goldbox}
\begin{minipage}[t]{0.48\linewidth}
    \textbf{Homework / Practice}
    \begin{itemize}
        \item Four scenarios: students check BINS and, where satisfied, compute
              a specified $P(X = k)$ showing formula setup.
        \item One problem asks students to compute $\binom{6}{2}$ by hand and
              explain what the number represents.
        \item AP-style free-response: a basketball shooting scenario with full
              BINS justification.
    \end{itemize}
\end{minipage}
\hfill
\begin{minipage}[t]{0.48\linewidth}
    \textbf{Extension / Preview}
    \begin{itemize}
        \item \textit{Extension}: Build the complete distribution table for
              $X \sim B(5, 0.4)$ and sketch its probability histogram.
        \item \textit{Preview}: Lesson 4.11 introduces the binomial mean
              $\mu_X = np$ and standard deviation $\sigma_X = \sqrt{np(1-p)}$
              --- start thinking about what the ``average number of successes''
              would be in today's quiz scenario.
    \end{itemize}
\end{minipage}
\end{skillbox}

\end{document}
